%%% FIELD proof-obstacle-compatibility
By the declared support types, \(X_0\) is compact and \(H\) is a nonempty compact subset of \(I\times(I-I)\).  Moreover, if \((d,z)\in H\), then \(d\in X_0\).  Indeed, let \(U\) be any relative open neighborhood of \(d\) in \(I\), and let \(V\) be any relative open neighborhood of \(z\) in \(I-I\).  Since \((d,z)\in\operatorname{supp}_P(X,Z)\),
\[
P\{(X,Z)\in U\times V\}>0.
\]
Therefore \(P(X\in U)>0\).  This holds for every such \(U\), so \(d\in\operatorname{supp}_P(X)=X_0\).

By \(\mathrm{ass:continuous\mbox{-}regression}\), \(g\) is continuous on the compact set \(H\), so there is a finite constant \(G_0\) such that \(|g(d,z)|\le G_0\) for every \((d,z)\in H\).  Compactness of \(I\) also bounds \(|z|\) and \(|x-d|\) on \(H\times X_0\).  Consequently, the supremum and infimum in \(\mathrm{def:obstacle\mbox{-}envelopes}\) are finite at every \(x\in X_0\).

For each \((d,z)\in H\), introduce the functions
\[
a_{d,z}(x)=g(d,z)-K|z|-L|x-d|,
\qquad
c_{d,z}(x)=g(d,z)+K|z|+L|x-d|.
\]
For all \(x,x'\in X_0\), the reverse triangle inequality gives
\[
|a_{d,z}(x)-a_{d,z}(x')|
=L\bigl||x-d|-|x'-d|\bigr|
\le L|x-x'|,
\]
and the same bound holds for \(c_{d,z}\).  If every member \(f_j\) of a pointwise finite family is \(L\)-Lipschitz, then
\[
\sup_j f_j(x)-\sup_j f_j(x')
\le \sup_j\{f_j(x)-f_j(x')\}
\le L|x-x'|.
\]
Interchanging \(x\) and \(x'\) proves that \(\sup_j f_j\) is \(L\)-Lipschitz.  Applying the same argument to \(-f_j\) proves the claim for a pointwise infimum.  Finally, the scalar maps \(v\mapsto\max\{-M,v\}\) and \(v\mapsto\min\{M,v\}\) are one-Lipschitz.  Therefore \(\ell_P\) and \(r_P\), as defined in \(\mathrm{def:obstacle\mbox{-}envelopes}\), are finite-valued and \(L\)-Lipschitz on \(X_0\), and in particular belong to \(C(X_0)\).

We now prove the asserted class equality.  Let \(u\in\mathcal C_P\).  Fix \(x\in X_0\) and \((d,z)\in H\).  Since \(d\in X_0\), the Lipschitz and support-cone clauses of \(\mathrm{def:feasible\mbox{-}trend\mbox{-}class}\) give
\[
u(x)
\ge u(d)-L|x-d|
\ge g(d,z)-K|z|-L|x-d|
=a_{d,z}(x).
\]
The level clause gives \(u(x)\ge -M\).  Taking the supremum over \((d,z)\in H\) and then the maximum with \(-M\) yields
\[
u(x)\ge \ell_P(x).
\]
Likewise,
\[
u(x)
\le u(d)+L|x-d|
\le g(d,z)+K|z|+L|x-d|
=c_{d,z}(x),
\]
and \(u(x)\le M\).  Taking the infimum over \((d,z)\in H\) and then the minimum with \(M\) yields
\[
u(x)\le r_P(x).
\]
Because \(x\) was arbitrary and membership in \(\mathcal C_P\) already imposes continuity and Lipschitz constant at most \(L\), we obtain
\[
\mathcal C_P
\subseteq
\{u\in C(X_0):\operatorname{Lip}(u)\le L,\ \ell_P\le u\le r_P\}.
\]

Conversely, let \(u\in C(X_0)\) satisfy \(\operatorname{Lip}(u)\le L\) and \(\ell_P\le u\le r_P\), where the inequalities are pointwise.  The outer truncations in \(\mathrm{def:obstacle\mbox{-}envelopes}\) imply
\[
\ell_P(x)\ge -M,
\qquad
r_P(x)\le M
\qquad (x\in X_0).
\]
Hence \(-M\le u(x)\le M\) for every \(x\in X_0\), so \(\|u\|_\infty\le M\).  For an arbitrary \((d,z)\in H\), evaluating the defining supremum and infimum at that same support point gives
\[
\ell_P(d)
\ge g(d,z)-K|z|-L|d-d|
=g(d,z)-K|z|
\]
and
\[
r_P(d)
\le g(d,z)+K|z|+L|d-d|
=g(d,z)+K|z|.
\]
Combining these inequalities with \(\ell_P(d)\le u(d)\le r_P(d)\) yields
\[
g(d,z)-K|z|\le u(d)\le g(d,z)+K|z|,
\]
or equivalently
\[
|g(d,z)-u(d)|\le K|z|.
\]
Thus \(u\) satisfies the level, Lipschitz, and support-cone clauses of \(\mathrm{def:feasible\mbox{-}trend\mbox{-}class}\), and therefore \(u\in\mathcal C_P\).  This proves
\[
\mathcal C_P
=
\{u\in C(X_0):\operatorname{Lip}(u)\le L,\ \ell_P\le u\le r_P\}.
\]

If \(\mathcal C_P\ne\varnothing\), choose \(u\in\mathcal C_P\).  The class equality just proved gives
\[
\ell_P(x)\le u(x)\le r_P(x)
\qquad\text{for every }x\in X_0,
\]
so \(\ell_P\le r_P\) pointwise.  Conversely, suppose \(\ell_P(x)\le r_P(x)\) for every \(x\in X_0\).  Set \(u=\ell_P\).  The first part of the proof shows that \(u\in C(X_0)\) and \(\operatorname{Lip}(u)\le L\); the assumed order gives \(\ell_P\le u\le r_P\).  The established class equality then implies \(u\in\mathcal C_P\).  Hence \(\mathcal C_P\ne\varnothing\), completing both directions of the compatibility criterion.
